% TOP_PROBLEM: {"id": 30006848, "title": "The Dynamic Optimality Conjecture for Splay Trees", "category_id": 15, "category": {"id": 15, "name": "computer_science", "display_name": "Computer Science", "description": "Computational complexity, algorithms, and theoretical CS.", "slug": "computer-science", "order_index": 15, "created_at": "2026-07-31T15:26:25.671Z"}, "status": "open"}
% FIRST_POSTED: 2026-09-25
% ENTRY_KIND: statement_only
% !TeX program = lualatex
% Definition and sources only; this file is not an LLM solution attempt.
\documentclass[11pt]{article}
\usepackage[margin=25mm]{geometry}
\usepackage{fontspec}
\setmainfont{Latin Modern Roman}
\usepackage{amsmath,amssymb,mathtools}
\usepackage{unicode-math}
\setmathfont{Latin Modern Math}
\usepackage{xurl,hyperref}
\hypersetup{colorlinks=true,urlcolor=blue}
\setcounter{secnumdepth}{0}
\setlength{\parindent}{0pt}
\setlength{\parskip}{5pt}
\setlength{\emergencystretch}{3em}
\begin{document}
\section*{\texorpdfstring{The Dynamic Optimality Conjecture for Splay Trees}{The Dynamic Optimality Conjecture for Splay Trees}}\label{30006848}
\subsection{Definitions and mathematical statement}
In the dynamic binary-search-tree model, a fixed ordered set of $n$ keys is stored in a rooted binary search tree. Servicing an access permits pointer moves and rotations, each charged unit cost under the source's conventions. Splaying moves the accessed key to the root by the standard zig, zig--zig, and zig--zag steps. For an access sequence $\sigma$, let $\operatorname{OPT}(\sigma)$ be the minimum cost over offline BST algorithms, which may know the complete sequence in advance. Dynamic optimality asks for an absolute constant $C$ such that
\[
 \operatorname{cost}_{\rm splay}(\sigma)
 \le C\operatorname{OPT}(\sigma)+I(n)
\]
for every $\sigma$, where $I(n)$ is the standard initialization allowance in the chosen model, independent of the access sequence. The input catalogue does not fix a numeric expression for $I$; a precise comparison must retain the cited source's initial-tree convention, rather than allowing an arbitrary sequence-dependent additive error.
\par\smallskip\textit{Formulation and concept references:} \href{https://arxiv.org/abs/2607.18498}{[S1]}.

\subsection{Short English statement}
Are splay trees always within a universal constant factor of the best binary-search-tree strategy that already knows the entire access sequence?
\subsection{Sources}
\begin{itemize}
\item[S1]Petr Chmel et al., Splay trees are almost dynamically optimal (2026).
\url{https://arxiv.org/abs/2607.18498}.
\textit{Formulation source.}
\end{itemize}
\end{document}
